% ============================================================
% CAPÍTULO 3 - TRABALHOS RELACIONADOS
% ============================================================
\chapter{Trabalhos Relacionados}
\label{cap:trabalhos}

A análise comparativa de arquiteturas de redes neurais para tarefas de visão computacional tem
sido um campo fértil de pesquisa, buscando estabelecer os modelos mais eficientes para
diferentes contextos de aplicação. Nesta seção, são apresentados trabalhos que, assim como
este, se dedicaram a comparar o desempenho de diferentes arquiteturas, seja em problemas de
classificação de imagens ou na detecção de objetos. A análise destes estudos é fundamental
para contextualizar a presente pesquisa, identificar as metodologias consolidadas e destacar
as lacunas que este trabalho se propõe a preencher.

% ------------------------------------------------------------
% 3.1 Análise Comparativa de CNNs para Triagem de Lesões de Pele
% ------------------------------------------------------------
\section{Análise Comparativa de CNNs para Triagem de Lesões de Pele}

Em um estudo focado na área médica, \citeonline{santos2019analise} realizou uma análise comparativa
entre as arquiteturas ResNet, VGG e GoogLeNet para a tarefa de triagem e classificação de
lesões de pele. O objetivo era avaliar a viabilidade de utilizar estas redes para apoiar o
diagnóstico médico a partir de um vasto conjunto de dados de 140.446 imagens obtidas em um
cenário real do sistema de saúde.

Após conduzir 42 cenários experimentais, variando arquiteturas, estratégias de inicialização
de pesos e taxas de aprendizado, o autor concluiu que a arquitetura GoogLeNet, com
inicialização randômica dos pesos, apresentou o melhor desempenho. O modelo atingiu 89,72\%
de acurácia para sugestão diagnóstica e 96,03\% para o encaminhamento adequado do paciente.

O trabalho de \citeonline{santos2019analise} é relevante por comparar arquiteturas de base
(\textit{backbones}) que também são utilizadas nos modelos de detecção de objetos, embora seu
foco seja estritamente a classificação.

% ------------------------------------------------------------
% 3.2 Estudo Comparativo de CNNs para Classificação de Imagens
% ------------------------------------------------------------
\section[Estudo Comparativo de CNNs para Classificação de Imagens]{Estudo Comparativo de CNNs para\\Classificação de Imagens}

De forma semelhante, \citeonline{silva2018estudo} conduziu um estudo comparativo direto entre as
arquiteturas VGG-16, ResNet-50 e Inception, aplicando-as a um problema de classificação de
imagens de um domínio específico: a identificação de personagens do desenho ``Os Simpsons''.
O autor buscou avaliar o poder de generalização dessas redes, pré-treinadas com o
\textit{dataset} ImageNet, em um conjunto de dados visualmente distinto.

Utilizando métricas como acurácia e \textit{log loss}, e testando diferentes otimizadores, o
estudo concluiu que a ResNet-50 obteve a maior acurácia (93,4\%), enquanto a Inception
alcançou o menor \textit{log loss} (0.26). Uma das principais constatações foi que a escolha
do otimizador (SGD em vez de RMSprop) foi crucial para o bom desempenho das redes. Este
trabalho corrobora a necessidade de uma análise empírica para determinar a arquitetura mais
adequada para um dado problema de classificação.

% ------------------------------------------------------------
% 3.3 Análise Comparativa de Modelos de Detecção de Objetos
% ------------------------------------------------------------
\section[Análise Comparativa de Modelos de Detecção de Objetos]{Análise Comparativa de Modelos de\\Detecção de Objetos}

Diferentemente dos trabalhos anteriores que focam na classificação,
\citeonline{su2025comparative} apresenta uma análise comparativa que aborda diretamente os três
modelos centrais deste TCC: Faster R-CNN, SSD e YOLO. O autor realiza uma revisão sistemática,
comparando as arquiteturas com base em dados de literatura sobre três eixos principais:
precisão (mAP), velocidade (FPS) e consumo de recursos (GPU VRAM).

A análise de \citeonline{su2025comparative} destaca o claro \textit{trade-off} entre os modelos: a
Faster R-CNN, uma abordagem de dois estágios, alcança alta precisão (73.2\% mAP) sendo ideal
para tarefas que não exigem tempo real, mas consome mais recursos; a YOLO, por sua vez,
prioriza a velocidade (45 FPS), sendo adequada para aplicações em tempo real, ao custo de uma
menor precisão (63.4\% mAP); por fim, a SSD se posiciona como uma solução de equilíbrio entre
velocidade e acurácia. Este artigo fornece um referencial teórico robusto para a comparação
que será realizada de forma prática neste trabalho.

% ------------------------------------------------------------
% 3.4 Aplicação de Redes Neurais na Identificação de Áreas Cafeeiras
% ------------------------------------------------------------
\section[Aplicação de Redes Neurais na Identificação de Áreas Cafeeiras]{Aplicação de Redes Neurais na\\Identificação de Áreas Cafeeiras}

A dissertação de \citeonline{andrade2011redes} demonstra uma aplicação prática de Redes Neurais
Artificiais (RNA) do tipo \textit{Multi-Layer Perceptron} (MLP) para um problema complexo de
classificação de imagens de satélite, com foco na identificação de áreas de cafeicultura. O
principal desafio era a alta similaridade espectral entre a cultura do café e a vegetação
nativa, o que dificultava a distinção por métodos tradicionais.

Utilizando uma metodologia que envolvia pré-processamento para mascarar áreas irrelevantes e
treinamento supervisionado com amostras representativas, o estudo alcançou um índice Kappa
considerado ``bom'' (entre 67\% e 71,85\%) para as regiões analisadas, superando outros
classificadores automáticos da literatura.

Embora utilize uma arquitetura mais antiga, o trabalho de \citeonline{andrade2011redes} é pertinente
por ilustrar a aplicação de redes neurais para resolver problemas de classificação com alta
confusão entre classes, um desafio também presente na detecção de múltiplos objetos em cenas
complexas.

% ------------------------------------------------------------
% 3.5 Síntese e Contribuição do Presente Trabalho
% ------------------------------------------------------------
\section[Síntese e Contribuição do Presente Trabalho]{Síntese e Contribuição do Presente\\Trabalho}

Em síntese, a literatura apresenta diversas análises comparativas que validam a eficácia de
diferentes arquiteturas de redes neurais para tarefas de visão computacional. Enquanto
trabalhos como os de \citeonline{santos2019analise} e \citeonline{silva2018estudo} se concentram na
comparação de arquiteturas para classificação de imagens, e \citeonline{andrade2011redes} demonstra
uma aplicação específica de redes neurais, o estudo de \citeonline{su2025comparative} oferece um
comparativo teórico direto dos modelos de detecção de objetos.

O presente TCC se diferencia ao realizar uma validação experimental e quantitativa dessas três
arquiteturas (YOLO, Faster R-CNN e SSD), aplicando-as aos principais \textit{datasets} de
referência (PASCAL VOC e COCO) e analisando criticamente seu desempenho não apenas em métricas
de precisão, mas também em tempo de inferência, preenchendo uma lacuna ao fornecer subsídios
práticos para a escolha do modelo mais adequado a diferentes cenários.
